% GENERATED FILE. Edit canonical lessons, then run npm run generate:books.
% canonical-source-hash: fnv1a64:2efd71f2d86a1af4
% canonical-lessons: RU-C96-towel, RU-C96-soap, RU-C96-toothbrush, RU-C96-comb, RU-C96-rubbish

\chapter{Towel, Soap, Toothbrush, Comb, Rubbish}
\label{ch:ru-towel-soap-toothbrush-comb-rubbish}
\hlchaptermodality{96}

\begin{chapteropening}
\textbf{By the end of this chapter:} \emph{I can say a towel, soap, a toothbrush, a comb and rubbish.}

Complete the last lesson of chapter 96: I can say a towel, soap, a toothbrush, a comb and rubbish.
\end{chapteropening}

\section[\texorpdfstring{\ru{полотенце} (poloténtse)}{полотенце (poloténtse)}]{\ru{полотенце} (poloténtse) — a towel}

\label{lesson:RU-C96-towel}



\begin{quote}
Before the new one: say the Russian for a pillow, then the Russian for a blanket.
\end{quote}

\subsection*{You'll want to know: \ru{полотенце}}
\textbf{\ru{полотенце}} (\emph{poloténtse}) — \textquotedblleft{}a towel\textquotedblright{}.

\begin{cousinweb}
\textbf{\ru{полотенце}} is \emph{\ru{полотно}}, \textquotedblleft{}linen cloth\textquotedblright{}, with the diminutive \emph{-\ru{енце}}. \emph{\ru{полотно}} $\leftarrow$ Proto-Slavic \emph{*poltĭno}, which Vasmer relates to Sanskrit \emph{paṭa-}, \textquotedblleft{}cloth\textquotedblright{}.
\end{cousinweb}

\begin{grammarlens}[title={What you've built}]
One more word for this chapter.
\end{grammarlens}

\subsection*{Your turn}
\begin{itemize}
\raggedright
  \item \emph{Say it:} \emph{poloténtse}
  \item \emph{Say it:} \emph{poloténtse}, once more
  \item \emph{Say it:} say \emph{poloténtse}
  \item \emph{Recall:} say the Russian for a pillow, then the Russian for a blanket, then say \emph{poloténtse} again
\end{itemize}

\begin{tcolorbox}[breakable,colback=teal!4,colframe=teal!35!black,title={Before you move on}]
What does \ru{полотенце} mean? (\textquotedblleft{}A towel\textquotedblright{}.) Say it once more.
\end{tcolorbox}
\section[\texorpdfstring{\ru{мыло} (mýlo)}{мыло (mýlo)}]{\ru{мыло} (mýlo) — soap}

\label{lesson:RU-C96-soap}



\begin{quote}
Before the new one: say the Russian for a blanket, then the Russian for a towel.
\end{quote}

\subsection*{You'll want to know: \ru{мыло}}
\textbf{\ru{мыло}} (\emph{mýlo}) — \textquotedblleft{}soap\textquotedblright{}.

\begin{cousinweb}
\textbf{\ru{мыло}} $\leftarrow$ Proto-Slavic \emph{*mydlo}, built on \emph{\ru{мыть}}, \textquotedblleft{}to wash\textquotedblright{} (Vasmer): the thing you wash with.
\end{cousinweb}

\begin{grammarlens}[title={What you've built}]
One more word for this chapter.
\end{grammarlens}

\subsection*{Your turn}
\begin{itemize}
\raggedright
  \item \emph{Say it:} \emph{mýlo}
  \item \emph{Say it:} \emph{mýlo}, once more
  \item \emph{Say it:} say \emph{mýlo}
  \item \emph{Recall:} say the Russian for a blanket, then the Russian for a towel, then say \emph{mýlo} again
\end{itemize}

\begin{tcolorbox}[breakable,colback=teal!4,colframe=teal!35!black,title={Before you move on}]
What does \ru{мыло} mean? (\textquotedblleft{}Soap\textquotedblright{}.) Say it once more.
\end{tcolorbox}
\section[\texorpdfstring{\ru{зубная} \ru{щётка} (zubnáya shchyótka)}{зубная щётка (zubnáya shchyótka)}]{\ru{зубная} \ru{щётка} (zubnáya shchyótka) — a toothbrush}

\label{lesson:RU-C96-toothbrush}



\begin{quote}
Before the new one: say the Russian for a towel, then the Russian for soap.
\end{quote}

\subsection*{You'll want to know: \ru{зубная} \ru{щётка}}
\textbf{\ru{зубная} \ru{щётка}} (\emph{zubnáya shchyótka}) — \textquotedblleft{}a toothbrush\textquotedblright{}.

\begin{cousinweb}
\textbf{\ru{зубная} \ru{щётка}} is \textquotedblleft{}tooth brush\textquotedblright{}: \emph{\ru{зубная}}, from \emph{\ru{зуб}}, \textquotedblleft{}tooth\textquotedblright{}, with \emph{\ru{щётка}}, \textquotedblleft{}brush\textquotedblright{}, from \emph{\ru{щеть}}, \textquotedblleft{}bristle\textquotedblright{}, the root of \emph{\ru{щетина}}, \textquotedblleft{}stubble\textquotedblright{}. Vasmer finds the proposed origins of that root doubtful.
\end{cousinweb}

\begin{grammarlens}[title={What you've built}]
One more word for this chapter.
\end{grammarlens}

\subsection*{Your turn}
\begin{itemize}
\raggedright
  \item \emph{Say it:} \emph{zubnáya shchyótka}
  \item \emph{Say it:} \emph{zubnáya shchyótka}, once more
  \item \emph{Say it:} say \emph{zubnáya shchyótka}
  \item \emph{Recall:} say the Russian for a towel, then the Russian for soap, then say \emph{zubnáya shchyótka} again
\end{itemize}

\begin{tcolorbox}[breakable,colback=teal!4,colframe=teal!35!black,title={Before you move on}]
What does \ru{зубная} \ru{щётка} mean? (\textquotedblleft{}A toothbrush\textquotedblright{}.) Say it once more.
\end{tcolorbox}
\section[\texorpdfstring{\ru{расчёска} (raschyóska)}{расчёска (raschyóska)}]{\ru{расчёска} (raschyóska) — a comb}

\label{lesson:RU-C96-comb}



\begin{quote}
Before the new one: say the Russian for soap, then the Russian for a toothbrush.
\end{quote}

\subsection*{You'll want to know: \ru{расчёска}}
\textbf{\ru{расчёска}} (\emph{raschyóska}) — \textquotedblleft{}a comb\textquotedblright{}.

\begin{cousinweb}
\textbf{\ru{расчёска}} is \emph{\ru{рас}-}, \textquotedblleft{}apart\textquotedblright{}, with \emph{\ru{чесать}}, \textquotedblleft{}to comb\textquotedblright{}, and \emph{-\ru{ка}}. \emph{\ru{чесать}} $\leftarrow$ Proto-Slavic \emph{*česati}, a cousin of Lithuanian \emph{kasýti}, \textquotedblleft{}to scratch\textquotedblright{} (Vasmer).
\end{cousinweb}

\begin{grammarlens}[title={What you've built}]
One more word for this chapter.
\end{grammarlens}

\subsection*{Your turn}
\begin{itemize}
\raggedright
  \item \emph{Say it:} \emph{raschyóska}
  \item \emph{Say it:} \emph{raschyóska}, once more
  \item \emph{Say it:} say \emph{raschyóska}
  \item \emph{Recall:} say the Russian for soap, then the Russian for a toothbrush, then say \emph{raschyóska} again
\end{itemize}

\begin{tcolorbox}[breakable,colback=teal!4,colframe=teal!35!black,title={Before you move on}]
What does \ru{расчёска} mean? (\textquotedblleft{}A comb\textquotedblright{}.) Say it once more.
\end{tcolorbox}
\section[\texorpdfstring{\ru{мусор} (músor)}{мусор (músor)}]{\ru{мусор} (músor) — rubbish}

\label{lesson:RU-C96-rubbish}



\begin{quote}
Before the new one: say the Russian for a toothbrush, then the Russian for a comb.
\end{quote}

\subsection*{You'll want to know: \ru{мусор}}
\textbf{\ru{мусор}} (\emph{músor}) — \textquotedblleft{}rubbish\textquotedblright{}.

\begin{cousinweb}
\textbf{\ru{мусор}} — Vasmer derives it from \emph{*mouds-} and compares Greek \emph{mýsos}, \textquotedblleft{}defilement\textquotedblright{}; Trubachev, in a note, ties it instead to dialect \emph{\ru{бусор}}, \textquotedblleft{}junk\textquotedblright{}.
\end{cousinweb}

\begin{grammarlens}[title={What you've built}]
One more word for this chapter.
\end{grammarlens}

\subsection*{Your turn}
\begin{itemize}
\raggedright
  \item \emph{Say it:} \emph{músor}
  \item \emph{Say it:} \emph{músor}, once more
  \item \emph{Say it:} say \emph{músor}
  \item \emph{Recall:} say the Russian for a toothbrush, then the Russian for a comb, then say \emph{músor} again
\end{itemize}

\begin{tcolorbox}[breakable,colback=teal!4,colframe=teal!35!black,title={Before you move on}]
What does \ru{мусор} mean? (\textquotedblleft{}Rubbish\textquotedblright{}.) Say it once more.
\end{tcolorbox}
